% !TeX root = ../main.tex
\chapter{Deployment Manifests and Simulator Specification}
\label{app:deployment}

This appendix collects the deployment files, the Telegraf bridge configuration, the QuestDB table definitions, and the specification used to extend the vehicle simulator. Chapter~\ref{chap:system_implementation} describes what those files do. The listings are the files of the prototype, copied so the manuscript can be read without the repository open beside it. Credentials in them are the local development defaults.

\section{Docker Compose}
\label{sec:app_compose}

The root file includes one Compose file per subsystem, declares the bridge networks, and mounts the partner certificate. The simulator include also applies the override that builds the image from source. Telegraf is the telemetry-bridge file: it is no longer a service of the MQTT broker. Training is the \texttt{ml} profile of the machine-learning module and does not start with \texttt{docker compose up}.

\lstinputlisting[language=yaml,caption={Root Compose file},label={lst:main-compose},basicstyle=\ttfamily\footnotesize,breaklines=true]{appendices/compose/root.yml}

\lstinputlisting[language=yaml,caption={MQTT broker},label={lst:mqtt-broker-compose},basicstyle=\ttfamily\footnotesize,breaklines=true]{appendices/compose/mqtt-broker.yml}

\lstinputlisting[language=yaml,caption={Vehicle simulator},label={lst:vda5050-simulator-docker-compose},basicstyle=\ttfamily\footnotesize,breaklines=true]{appendices/compose/simulator.yml}

\lstinputlisting[language=yaml,caption={Simulator build override},label={lst:simulator-override},basicstyle=\ttfamily\footnotesize,breaklines=true]{appendices/compose/simulator-override.yml}

\lstinputlisting[language=yaml,caption={Redpanda},label={lst:redpanda-docker-compose},basicstyle=\ttfamily\footnotesize,breaklines=true]{appendices/compose/redpanda.yml}

\lstinputlisting[language=yaml,caption={Spark streaming jobs},label={lst:spark-docker-compose},basicstyle=\ttfamily\footnotesize,breaklines=true]{appendices/compose/spark.yml}

\lstinputlisting[language=yaml,caption={MinIO},label={lst:minio-docker-compose},basicstyle=\ttfamily\footnotesize,breaklines=true]{appendices/compose/minio.yml}

\lstinputlisting[language=yaml,caption={Telemetry bridge},label={lst:telemetry-bridge-compose},basicstyle=\ttfamily\footnotesize,breaklines=true]{appendices/compose/telemetry-bridge.yml}

\lstinputlisting[language=yaml,caption={QuestDB},label={lst:questdb-compose},basicstyle=\ttfamily\footnotesize,breaklines=true]{appendices/compose/questdb.yml}

\lstinputlisting[language=yaml,caption={Grafana},label={lst:grafana-compose},basicstyle=\ttfamily\footnotesize,breaklines=true]{appendices/compose/grafana.yml}

\lstinputlisting[language=yaml,caption={MLflow},label={lst:mlflow-compose},basicstyle=\ttfamily\footnotesize,breaklines=true]{appendices/compose/mlflow.yml}

\lstinputlisting[language=yaml,caption={Training module},label={lst:ml-module-compose},basicstyle=\ttfamily\footnotesize,breaklines=true]{appendices/compose/ml-module.yml}

\lstinputlisting[language=yaml,caption={Factory orchestrator},label={lst:orchestrator-compose},basicstyle=\ttfamily\footnotesize,breaklines=true]{appendices/compose/factory-orchestrator.yml}

\lstinputlisting[language=yaml,caption={Digital twin},label={lst:twin-compose},basicstyle=\ttfamily\footnotesize,breaklines=true]{appendices/compose/digital-twin.yml}

\lstinputlisting[language=yaml,caption={Platform portal},label={lst:portal-compose},basicstyle=\ttfamily\footnotesize,breaklines=true]{appendices/compose/platform-portal.yml}

\section{Telegraf configuration}
\label{sec:app_telegraf}

The bridge in Section~\ref{subsec:telegraf_integration} is configured by this file. It is mounted into the Telegraf service of Listing~\ref{lst:telemetry-bridge-compose}.

\lstinputlisting[language=configToml,caption={Telegraf bridging configuration (telegraf.conf)},label={lst:telegraf-conf},basicstyle=\ttfamily\footnotesize,breaklines=true]{appendices/config/telegraf.conf}

\section{QuestDB tables}
\label{sec:app_questdb_schema}

\texttt{questdb-init} applies these two files, as described in Section~\ref{subsec:questdb-deployment-schema}.

\lstinputlisting[language=SQL,caption={QuestDB table \texttt{fleet\_features}},label={lst:fleet-features-sql},basicstyle=\ttfamily\footnotesize,breaklines=true]{appendices/schema/fleet_features.sql}

\lstinputlisting[language=SQL,caption={QuestDB table \texttt{fleet\_anomalies}},label={lst:fleet-anomalies-sql},basicstyle=\ttfamily\footnotesize,breaklines=true]{appendices/schema/fleet_anomalies.sql}

\section{Specification used to extend the simulator}
\label{app:simulator_prompt}

The electromechanical extension of the simulator was specified to the model as the document below, and the generated code was reviewed against it before it was kept. The first draft asked for a separate physics task, a full armature inductance, a polarisation branch in the battery, and a stator short. The integrated simulator does not do those, and it adds wheel wear as a drop in rolling resistance, bearing wear, a housing-vibration RMS, the accelerated discharge, and the depletion stop. Those departures are marked in the document as implemented deviations, so the text matches the simulator that publishes the telemetry. The equations used in Chapter~\ref{chap:system_implementation} are the ones this specification settles on.

\lstinputlisting[caption={Simulator extension specification},label={lst:simulator-prompt},basicstyle=\ttfamily\footnotesize,breaklines=true]{appendices/prompt/complete_sensor_adaptation.md}
